% !TEX root = ../monograph.tex
% Fragment for Chapter13: insert after the coarse separated bound.
\section{Sharp short-distance bounds for a fixed number of insertions}
\label{ch13:sharp-section}

For a fixed number of insertions the loss in
\eqref{ch13:power} can be reduced to an arbitrarily small excess over
the sum of their degrees. This uses the normalized integral over the reflected tiling,
not the stronger Gaussian comparison used for the separating families. We keep
all incoming joint components with insertions; replacing them by a product of
terms of source degree one would invalidate the estimate.

Write $\tau_m=\log(1/a_m)$, so
$\widehat\beta_m=\bar\beta+\kappa_G\tau_m$. Let $C_j\ge1$ include
the inverse-column, finite-section and source constants at stage $j$.
The stage ratio and the lower bound on the reference coupling may be chosen to satisfy
\begin{equation}
 \frac{\log C_j}{\log \mathscr L_j}\longrightarrow0,\qquad
 B_{i+1}^2\le2B_i^2,\qquad
 \frac{(L_iB_i^2)^4 q_{i+1}^{\rm site}}{\tau_{i+1}}
                          \longrightarrow0.
 \label{ch13:sharp-schedule}
\end{equation}
Here $B_i$ is the support-block radius, and $q_i^{\rm site}\ge1$ is
the per-site normalization exponent of \eqref{ch11:replicated-bound};
it is not a density with insertions or the coefficient of the coupling constant renormalization.
For the path-action representation of the density one can use
$q_i^{\rm site}\le C_{J(i)}(1+\log\widehat\beta_i)$.
Indeed the exponent without its constant term is at most $C_{J(i)}V$. Set
$\rho=\widehat\beta_i^{-1/2}$ and restrict every link to
$\|U_e-1\|_G\le\rho/4$, using the matrix norm \eqref{ch01:matrix-norm}.
On this event the exponent is at least $-C_{J(i)}V$ and all large-field
indicators vanish. Exponential summability in the path length bounds the
leading positive path quadratic by $C_{J(i)}\widehat\beta_i\rho^2V$.
Signed regular costs and the absolute
activity site sum satisfy the bounds \eqref{ch04:activity-bounds}. The lower bound for the partition function
of Section~\ref{ch11:replicated} therefore adds only
$C_{J(i)}V+2D_GV\log\widehat\beta_i$ to the exponent, where $D_G$ is an ambient
real matrix dimension for the gauge group.

These requirements are compatible with sufficiently large lower bounds on the reference coupling. At each fixed
stage $B_i=O_j((\log(2+\tau_i))^2)$ and
$q_i^{\rm site}=O_j(1+\log(2+\tau_i))$, so the last quotient in
\eqref{ch13:sharp-schedule} tends to zero at its fixed ratio. Choose
the ratio after $C_j$, then increase the lower bound on the reference coupling until that quotient is at
most $1/j$ for both adjacent-stage choices. The same choice gives
the displayed bound on consecutive $B_i$,
$\log L_i=o(\tau_i)$, and $\log C_{J(i+1)}=o(\tau_i)$
at transitions. Choose each ratio against both adjacent-stage column
constants, which are fixed beforehand. These are conditions on the single selected
schedule; no new schedule is chosen after fixing a correlation.

\subsection{A normalized local insertion}

Fix a scaling-dimension bound $d$ and a ball of radius $R$ in the
coefficient space of $V_d$. Set
\[
 E_k=C_{d,R}^{\rm sec} C_{J(k)}\prod_{i<k}C_{J(i+1)},\qquad
 \omega_k=a_k^d\widehat\beta_k^{-w_d}E_k^{-1}.
\]
The fixed section constant $C_{d,R}^{\rm sec}$ includes the coefficient radius and all its signed and
conjugated images. Choose the fixed integer $w_d$ large enough to absorb
the finite powers of the inverse small-field threshold in the degree-$d$
section and its counterpart with a prescribed large-field set, as well as the source attachment bound. Multiplication of the single-insertion inverse bounds gives the
factor $E_k$; the scaling normalization is exactly $a_k^{-d}$.
Thus multiplication of the \emph{full} insertion by $\omega_k$
makes the exponential bound for its regular term of source degree one
less than three and the bound for its large-field term of source degree one
less than one. Its remaining normalized size can be placed below the
attachment bound by increasing $w_d$.

The bounds in Chapter~\ref{ch11:insertions} for terms of source degree at least two retain every
incoming source term and every original eligibility event. Decreasing all source
weights only decreases their sums of absolute coefficients. The source attachment
inequalities at the intervening scales propagate this chosen small
weight; they require the same inverse-column and threshold powers already
included in $E_k,w_d$. Consequently \eqref{ch11:root-bound}, with
$x=3,y=1$, gives
\begin{equation}
 \|O_{N,p}^{(d)}(c)\|_{\mathrm{cell}(m)}
 \le \omega_k^{-1}
           \exp\bigl((a_m/a_k)^4q_k^{\rm site}+6\bigr).
 \label{ch13:sharp-root}
\end{equation}
It holds for all sufficiently large $N$ uniformly in period and in the coefficient
ball, whenever the reflected insertion points at the stopping scale have spacing
at least $B_k^2$. Their spacing stays adequate at finer scales because
$L_i\ge3$ whereas $B_{i+1}^2/B_i^2\le2$.

The logarithm of $E_k$ is $o(\tau_k)$: split its sum into a finite
initial segment and a tail where
$\log C_{J(i+1)}\le\epsilon\log L_i$, then let $\epsilon\downarrow0$.
The separate terminal column is absorbed by the same ratio bounds and
$\log L_{k-1}=o(\tau_k)$. Also
$\log\widehat\beta_k=o(\tau_k)$. Therefore
\begin{equation}
                        6-\log\omega_k=d\tau_k+o(\tau_k).
 \label{ch13:weight-size}
\end{equation}
This is the reason the prefactor from the reflected tiling does not introduce a fixed
extra scaling degree.

\subsection{The stopping scale and centered estimate}

Choose $k=k(m)$ as the first scale $k\ge m$ with
$B_k^2\le a_m/a_k$. It exists because $B_k$ grows subexponentially
in $\tau_k$, whereas $a_m/a_k=e^{\tau_k-\tau_m}$. By minimality,
\[
 (a_m/a_k)^4q_k^{\rm site}
       \le (L_{k-1}B_{k-1}^2)^4q_k^{\rm site}=o(\tau_k).
\]
If the first scale already qualifies the same conclusion follows from
the fixed-stage bound. Since $q_k^{\rm site}\ge1$,
\[
 0\le\tau_k-\tau_m=\log(a_m/a_k)
       \le(a_m/a_k)^4q_k^{\rm site}=o(\tau_k).
\]
It follows that $\tau_k/\tau_m\to1$.
Combining this with \eqref{ch13:sharp-root} and
\eqref{ch13:weight-size} proves the following statement.

\begin{proposition}[Sharp centered bound]
\label{ch13:sharp-centered}
For every scaling-dimension bound $d$, coefficient radius $R$ and $\epsilon>0$, there
is $C_{d,R,\epsilon}$, fixed before the cell scale $m$, such that
\[
 \|O_{N,p}^{(d)}(c,y)\|_{\mathrm{cell}(m,y)}
        \le C_{d,R,\epsilon}a_m^{-d-\epsilon},\qquad\|c\|\le R.
\]
For each fixed $m$ this holds for all sufficiently large cutoffs,
uniformly in $p,y$ and in signed and conjugated images. The insertion
is interior in that cell on the same tail.
\end{proposition}
\begin{proof}
For all sufficiently large $m$, the logarithm of the right side of
\eqref{ch13:sharp-root} is at most $(d+\epsilon)\tau_m$ by the
preceding limits. For the finite remaining set of $m$, use its chosen
fixed stopping scale and take the maximum of the finite constants.
The fixed link support becomes interior at every fixed cell
scale in continuum units as $N\to\infty$. Translation covariance makes the lower bound on $N$
independent of $y$.
\end{proof}

\subsection{Arbitrary separated points}

The slow ratio condition gives a sharper estimate of the first scale below a length
than \eqref{ch13:scale-thickness}. For any fixed $c>0$, let $m(r)$
be the first scale with $a_m<cr$. Then
\begin{equation}
            \frac{\log(1/a_{m(r)})}{\log(1/r)}\longrightarrow1
                         \quad(r\downarrow0).
 \label{ch13:first-below}
\end{equation}
Indeed $a_{m-1}\ge cr>a_m=a_{m-1}/L_{m-1}$, and the logarithm of
the overshoot is at most $\log L_{m-1}=o(\tau_{m-1})$.
Any fixed finite iteration of such choices of a first scale has the same property.

For a fixed number $n$ of insertions and mutually separated insertion points, first choose a
cell side below a fixed small multiple of their minimum separation $r$.
Translate its boundaries away from all insertion points. Then choose a centered
scale smaller than this side divided by $6(n+1)$.
The interior-cell comparison and the chessboard inequality apply.
Both chosen scales have logarithm $(1+o(1))\log(1/r)$ by
\eqref{ch13:first-below}. Apply Proposition~\ref{ch13:sharp-centered}
with a sufficiently small excess to each degree separately. We obtain,
for every $\epsilon>0$,
\begin{equation}
 |S_n(A_1,\ldots,A_n;x)|
    \le C_{\boldsymbol A,n,\epsilon}
             r^{-\sum_i d_i-\epsilon},\qquad 0<r\le1,
 \quad A_i\in F_{d_i},\quad\min_{i\ne j}|x_i-x_j|_\infty\ge r.
 \label{ch13:sharp-separated}
\end{equation}
The same bound is uniform over bounded sets of fixed raw representatives.
Its constant is independent of the diameter of the configuration.
For fixed derivatives, replace the differentiated field by its degree-raised
label; their total order is added to $\sum_i d_i$.

